\epanalipsi{Γ3}

\tmima{Λεξικό του κεφαλαίου}

\renewcommand{\arraystretch}{1.5}
\begin{longtable}{@{}>{\raggedright\arraybackslash}p{0.28\linewidth} >{\raggedright\arraybackslash}p{0.62\linewidth}@{}}
\toprule
\textbf{Λέξη ή σύμβολο} & \textbf{Με μία πρόταση}\\
\midrule
\endhead
πείραμα τύχης & Πείραμα που δεν μπορούμε να προβλέψουμε με βεβαιότητα.\\
δειγματικός χώρος & Το σύνολο όλων των δυνατών αποτελεσμάτων, μέσα σε άγκιστρα $\{\ \}$.\\
ενδεχόμενο & Ένα σύνολο στοιχείων του δειγματικού χώρου.\\
δενδροδιάγραμμα & Δέντρο με κλαδιά για κάθε επιλογή· κάθε μονοπάτι είναι ένα αποτέλεσμα.\\
πίνακας διπλής εισόδου & Πίνακας δύο σταδίων· κάθε κελί είναι ένα αποτέλεσμα.\\
ισοπίθανα & Αποτελέσματα με την ίδια πιθανότητα.\\
$\mathrm{Ρ}(\mathrm{Α})=\frac{\text{ευνοϊκά}}{\text{δυνατά}}$ & Κλασικός ορισμός της πιθανότητας (για ισοπίθανα).\\
βέβαιο, αδύνατο & Πιθανότητα $1$ ($100\%$) και πιθανότητα $0$.\\
σχετική συχνότητα & $\dfrac{\kappa}{\nu}$: πόσες φορές συνέβη κάτι διά το πλήθος των επαναλήψεων.\\
\bottomrule
\end{longtable}

\tmima{Ασκήσεις επανάληψης}

\ex Ποια είναι πειράματα τύχης; α) Τραβάω ένα φύλλο από ανακατεμένη τράπουλα. β) Υπολογίζω το $5+7$. γ) Κληρώνω το όνομα του παιδιού που θα κρατήσει τη σημαία. δ) Κοιτάζω την ημερομηνία στο ημερολόγιο.

\ex Επιλέγουμε τυχαία έναν μήνα του έτους που έχει $30$ ημέρες. Γράψε τον δειγματικό χώρο.

\ex Ρίχνουμε ένα νόμισμα και μετά ένα ζάρι. Σχεδίασε δενδροδιάγραμμα. Πόσα αποτελέσματα υπάρχουν; Πόσα από αυτά έχουν Κεφάλι και άρτιο αριθμό και ποια είναι η πιθανότητά τους;

\ex Σε μια σχολική γιορτή το μενού έχει $3$ πρώτα πιάτα, $4$ κυρίως και $2$ επιδόρπια. Πόσα διαφορετικά γεύματα μπορεί να φτιάξει ένας μαθητής με ένα πιάτο από κάθε κατηγορία;

\ex Ρίχνουμε δύο ζάρια. Βρες την πιθανότητα: α) το άθροισμα να είναι $10$, β) το άθροισμα να είναι τουλάχιστον $10$.

\ex Ένα κουτί έχει $5$ κόκκινες, $3$ μπλε και $2$ πράσινες μπάλες. Βρες την πιθανότητα να βγει: α) μπλε, β) όχι κόκκινη, γ) κίτρινη.

\ex Στην κλήρωση της γιορτής πουλήθηκαν $50$ λαχνοί και κερδίζουν οι $5$. Βρες την πιθανότητα να κερδίσει ένας λαχνός και να μην κερδίσει, ως κλάσμα, δεκαδικό και ποσοστό.

\ex Τραβάμε ένα φύλλο από μια τράπουλα $52$ φύλλων. Βρες την πιθανότητα να είναι σπαθί και να είναι εξάρι.

\ex Ένας τροχός έχει $10$ ίσα τμήματα αριθμημένα από το $1$ ως το $10$. Βρες την πιθανότητα να σταματήσει: α) σε πολλαπλάσιο του $3$, β) σε πρώτο αριθμό, γ) σε αριθμό μεγαλύτερο από $10$.

\ex Σε μια σακούλα με $30$ καραμέλες η πιθανότητα να βγει φράουλα είναι $0{,}4$ και να βγει λεμόνι $\frac{1}{3}$. Πόσες καραμέλες είναι φράουλα, πόσες λεμόνι και πόσες άλλες;

\ex Ρίχνουμε ένα νόμισμα τρεις φορές. Βρες την πιθανότητα να φέρουμε Γράμματα ακριβώς μία φορά και την πιθανότητα και οι τρεις ρίψεις να έχουν το ίδιο αποτέλεσμα.

\ex Ο Νίκος έριξε ένα ζάρι $300$ φορές και έφερε $6$ τις $52$ φορές. Βρες τη σχετική συχνότητα του $6$ και σύγκρινέ τη με την πιθανότητα.

\ex Η Ρήνη λέει ότι στις δύο ρίψεις νομίσματος η πιθανότητα να φέρουμε ένα Κεφάλι και ένα Γράμματα είναι $\frac{1}{3}$, γιατί τα αποτελέσματα είναι $\mathrm{ΚΚ}$, $\mathrm{ΚΓ}$, $\mathrm{ΓΓ}$. Βρες το λάθος και τη σωστή πιθανότητα.

\ex Με τα ψηφία $1,2,3$ φτιάχνουμε τριψήφιους αριθμούς με όλα τα ψηφία διαφορετικά. Σχεδίασε δενδροδιάγραμμα και βρες την πιθανότητα ένας τυχαίος από αυτούς να είναι άρτιος.

\ex Σωστό ή λάθος; α) Η πιθανότητα του αδύνατου ενδεχομένου είναι $0$. β) Τα αθροίσματα $2,3,\dots,12$ δύο ζαριών είναι ισοπίθανα. γ) Η πιθανότητα ενός ενδεχομένου είναι πάντα από $0$ ως $1$. δ) Αν $\mathrm{Ρ}(\mathrm{Α})=0{,}2$, τότε η πιθανότητα να μη συμβεί το $\mathrm{Α}$ είναι $0{,}8$. ε) Σε τέσσερις ρίψεις ενός νομίσματος υπάρχουν $8$ αποτελέσματα.

\ex Σωστό ή λάθος; α) Ένα δενδροδιάγραμμα με $3$ επιλογές στο πρώτο στάδιο και $2$ στο δεύτερο έχει $6$ τελικά αποτελέσματα. β) Η σχετική συχνότητα ενός ενδεχομένου είναι πάντα ίση με την πιθανότητά του. γ) Αν όλα τα αποτελέσματα είναι ευνοϊκά, η πιθανότητα είναι $1$. δ) Στις δύο ρίψεις νομίσματος το $\mathrm{ΚΓ}$ και το $\mathrm{ΓΚ}$ είναι διαφορετικά αποτελέσματα. ε) Αν ένα ενδεχόμενο έχει πιθανότητα $\frac{1}{2}$, σε $100$ επαναλήψεις θα συμβεί ακριβώς $50$ φορές.

\tmima{Χάρτης επανάληψης}

\begin{center}
\renewcommand{\arraystretch}{1.3}
\begin{tabular}{ll}
\toprule
\textbf{Ασκήσεις} & \textbf{Παράγραφος}\\
\midrule
1, 2, 13 & Γ3.1 Πειράματα τύχης και δειγματικός χώρος\\
3, 4, 14 & Γ3.2 Δενδροδιάγραμμα και πίνακας\\
5, 6, 7, 8, 9, 10, 11, 12 & Γ3.3 Κλασικός ορισμός της πιθανότητας\\
15, 16 & Όλες οι παράγραφοι του κεφαλαίου\\
\bottomrule
\end{tabular}
\end{center}
